% Consolidated formalization of 26 September 2026.
% Provenance and conditions: COMUN/procedencia_20260926/.
\clearpage
\section{Joint normalization of action, length, and gravitational coupling}
\label{md:rigidez}

The determination of a quantity is preserved when the readouts of the
same state are composed. Here the inscribed return, the action section,
the constitutive speed, and the positive route character are brought
together. Their common antecedent is APP--TRIT--TPK generation, with
its sheets, orientation, carry, memory, and prolongations; the regional
coordinates and $\alpha$ enter as outputs of that generation. The
following construction uses those outputs before performing the
dimensional recognition of $G$.

The normalization problem comprises three distinct operations: producing
a length from the recorded return, transporting the character without
losing its complement, and comparing the resulting energy and length
readouts. The identities in this section jointly fix those operations.
Their proofs concern finite fibres or bounded, uniformly positive
operators; the null sector retains its separate treatment.

\subsection{Inscribed return and longitudinal section}
\label{md:rigidez:longitud}

The marked carrier of the transport retains the support
\[
 \Omega=\{(j,k,q,u,v):j\in\mathbb Z/12\mathbb Z,\ 
 k-j\in\{0,3,6,9\},\ 1\le q\le8,\
 u,v\in\{1,2,4,5,7,8\},\ u<v\}.
\]
The index $j$ retains the dodecaphase phase; the four displacements
retain the quarter-turn, and the remaining indices retain the marked
incidence. Thus $N=|\Omega|=12\cdot4\cdot8\binom62=5760$ counts
incidence positions. The space of this register is transported together
with genealogical memory; its finite dimension belongs to this particular
carrier.

Let $B:V_U\to V_K$ be the transported Hadamard inverter,
$B^*B=BB^*=I/4$, $u_K=N^{-1/2}(1,\ldots,1)$,
$P_K=u_Ku_K^*$ and $Q_K=I-P_K$. Define
\[
 C_0=Q_KB,\qquad M_0=P_KB,\qquad U=2B.
\]
Orthogonality of the two projectors proves
\[
 C_0^*C_0+M_0^*M_0=I/4,\qquad C_0^*M_0=0,\qquad U^*U=UU^*=I.
\]
Thus $x\mapsto(2C_0x,2M_0x)$ preserves the complete archive of this
transport. $U$ is its unitary normalization; the enriched evolution
$U_t$ retains its own domain and update law.

Let $E_i=e_ie_i^*$ be the resolution of the register and
$\Delta_\Omega(A)=\sum_iE_iAE_i$ its diagonal expectation.

\begin{lemma}[Inscribed longitudinal return]
\label{md:rigidez:retorno}
The return and its inscribed readout satisfy
\[
 C_0C_0^*=Q_K/4,\qquad
 \Delta_\Omega(C_0C_0^*)=r_NI,\qquad
 r_N=\frac{N-1}{4N}=\frac{5759}{23040}.
\]
With the local scalar norm $n_H=\alpha^{16}$ of the cokernel of $H_4$,
the longitudinal composition on a length cochain $\beta_*$ is
\[
 \mathscr T_L(\beta_*)=
 n_H\bigl(\Delta_\Omega(C_0C_0^*)\otimes I\bigr)\beta_*
 =q_\alpha\beta_*,
 \qquad q_\alpha=\alpha^{16}r_N.
 \label{md:rigidez:qalpha}
\]
\end{lemma}
\begin{proof}
The equality $BB^*=I/4$ gives $C_0C_0^*=Q_K/4$. Every diagonal entry
of $P_K$ equals $1/N$, proving the expectation. Multiplying a cochain
by that scalar preserves its longitudinal type and gives the last
equality. If $\beta_*(\bar a)=-U_a\beta_*(a)$, the image satisfies
the same reversal. Likewise,
$\beta_m(a)=\sum_jU_j^{-1}\beta_{m+1}(a_j)$ yields the refined identity
by multiplying both sides by $q_\alpha$.
\end{proof}

The expectation has an explicit inscribed realization:
$We_i=e_i\otimes e_i$ and
$W^*(A\otimes I)W=\Delta_\Omega(A)$. For
$R_K=Q_K/4$ and $P_W=WW^*$, the complement
$Z_\perp=(I-P_W)(R_K\otimes I)W$ satisfies
\[
 (R_K\otimes I)W=Wr_N+Z_\perp,\qquad W^*Z_\perp=0,\qquad
 Z_\perp^*Z_\perp=\frac{N-1}{16N^2}I.
\]
Indeed, $R_K^2=R_K/4$, so the squared full norm is $r_N/4$,
whereas the squared norm of the inscribed component is $r_N^2$.
Their difference is precisely the stated complement. Preservation of
that complement accompanies the longitudinal readout; adding it to
another readout would change the observable.

The expectation is unique among bimodular maps into the diagonal algebra
that fix its projectors: on a matrix unit $E_{ij}$, bimodularity requires
$\mathcal E(E_{ij})=E_i\mathcal E(E_{ij})E_j$, which is zero if $i\ne j$,
and $\mathcal E(E_{ii})=E_{ii}$. This uniqueness preserves the marked
resolution.

With the reference longitudinal section $L_*$ already specified, one obtains
\begin{equation}
 L_\alpha=q_\alpha L_*,\qquad D=L_\alpha U.
 \label{md:rigidez:longitud-def}
\end{equation}
The readout operator uses the return as a linear operator on a length.
A quadratic readout of the source would produce a different type of
quantity. The section $L_*$, the inscribed resolution, and the local
norm are explicit antecedents of this composition.

\subsection{Polar transport and reciprocity of the same character}
\label{md:rigidez:polar}

Let $X$ be the operator character of the non-null positive sector,
retaining the route and its normalizations. In finite dimension $X>0$
is required; in infinite dimension $X$ must be bounded, self-adjoint,
and satisfy $X\ge\epsilon I$ for some $\epsilon>0$. Write $Y=UXU^*$.
The polar radial representation is
\[
 R=(DX^2D^*)^{1/2}.
\]
The positive square root is unique, and substituting $D=L_\alpha U$
yields $R=L_\alpha Y$. Equivalently,
$\|Rv\|^2=\|XD^*v\|^2$ for every $v$: polarization recovers
$R^2=DX^2D^*$, and positivity determines $R$. This condition specifies
the induced metric that is preserved.

The action $\hbar$ and the constitutive speed $c$ belong to the same
section. The compatible circular realization fixes
\begin{equation}
 \omega_0=\frac{c}{L_\alpha},\qquad
 t_0=\frac{\pi_{\rm HMT}L_\alpha}{54c},\qquad
 R_{108}=\frac{54}{\pi_{\rm HMT}}ct_0=L_\alpha.
 \label{md:rigidez:reloj}
\end{equation}
The length has been constructed before the circular clock is expressed.
This equality allows the notation $R_{108}$ to be retained in subsequent
developments and fixes its normalization. The dimensional bases are
transported jointly; $L_*$ and the edge basis $u_L$ have different roles.

\begin{theorem}[Joint determination of radial coupling]
\label{md:rigidez:teorema}
The readout operators on the same character,
\[
 H=\frac{\hbar c}{L_\alpha}Y,\qquad
 \Lambda=L_\alpha Y^{-1},\qquad R=L_\alpha Y,\qquad M=H/c^2,
\]
satisfy
\begin{equation}
 H\Lambda=\hbar cI,\qquad R\Lambda=L_\alpha^2I,\qquad
 R=\frac{L_\alpha^2}{\hbar c}H.
 \label{md:rigidez:identidades}
\end{equation}
In the dimensional identification of $R$ as the reduced gravitational
radius, $R=GM/c^2$, the common coefficient is unique:
\begin{equation}
 \boxed{G=\frac{c^3L_\alpha^2}{\hbar}.}
 \label{md:rigidez:G}
\end{equation}
\end{theorem}
\begin{proof}
Multiplication of the readout operators gives the first two identities.
Multiplying them on the right by $\Lambda^{-1}$ proves the third.
Substituting $M=H/c^2$ identifies the coefficient. If $G'$ preserves
the same readouts, $(G'-G)M/c^2=0$; in the non-null fibre,
invertibility of $M$ implies $G'=G$.
All products are defined on the complete fibre under the hypotheses
of the statement.
\end{proof}

The proof applies to every positive mode of this realization, without
choosing the electronic mass to normalize $G$. The identification of
the reduced radius is explicit; the radius $2GM/c^2$ corresponds to
another convention for reading the same constant. The null sector
remains outside the inverse. Extension to unbounded operators requires
a development of domains and limits distinct from the present theorem.

The square of $L_\alpha$ comes from the product of conjugate lengths.
Action enters inversely when $\Lambda$ is eliminated between energy
and length. The passage from energy to mass and from mass to radius
contributes $c^4$, which, combined with the $c$ in $H\Lambda=\hbar cI$,
produces $c^3$. The two normalizations fix the dimensionless factor
at one. Dimensional analysis verifies the exponents; composition
determines the coefficient.

\subsection{Planck scales and specialization of the unit character}
\label{md:rigidez:escalas}

The constructed length, the action section, and the speed determine
\[
 t_P=\frac{L_\alpha}{c},\qquad
 m_{P,a}=\frac{\hbar_a}{cL_\alpha},\qquad
 E_{P,a}=\frac{\hbar_a c}{L_\alpha}.
\]
Substituting \eqref{md:rigidez:G} into
$\sqrt{\hbar_aG_a/c^5}$, $\sqrt{\hbar_ac/G_a}$ and
$\sqrt{\hbar_ac^5/G_a}$ proves these identities, respectively;
positivity fixes each square root. The length precedes the conventional
recognition of these scales.
For a mode of character $y>0$,
\[
 m(y)=m_{P,a}y,\qquad r_g(y)=L_\alpha y,\qquad
 \bar\lambda(y)=L_\alpha/y.
\]
Equality of the two lengths is equivalent to $y=y^{-1}$ and hence
to $y=1$. The areal product is preserved for every positive mode;
coincidence of radii characterizes the unit mode.

The chronological step satisfies $t_0=(\pi_{\rm HMT}/54)t_P$;
the phase period is $108t_0=2\pi_{\rm HMT}t_P$ and
$\omega_{108}=1/t_P$. A horizon $R_h=\zeta_hL_\alpha$ retains
its own geometric condition. For the Schwarzschild specialization
of the same mass, $R_h=2r_g=2L_\alpha y$.
Comparing sections with fixed $L_\alpha,c$ and
$\rho_a=\hbar_b/\hbar_a$ yields
$G_b=G_a/\rho_a$, $m_{P,b}=\rho_a m_{P,a}$ and
$E_{P,b}=\rho_a E_{P,a}$. This comparison of sections preserves
$L_\alpha$ and $t_P$.

\subsection{Area and reconstruction of the memory sectors}
\label{md:rigidez:area}

Boundary incidence retains the grouping of six trits.
If $x_j=r_{2j-1}+3r_{2j}$, $r_i\in\{0,1,2\}$, the map
\[
 (x_1,x_2,x_3)\longmapsto
 (r_1+3r_2+9r_3,\ r_4+3r_5+9r_6)
\]
is a bijection between $(\mathbb Z/9\mathbb Z)^3$ and
$(\mathbb Z/27\mathbb Z)^2$. Its inverse extracts the six digits
in order, preserving the position at which carry is recorded. On the
boundary torus, the block $A=\{0,\ldots,8\}^2$ has two collections of links:
\begin{align*}
 \partial A_{90}&=\{((-1,j),(0,j)),((8,j),(9,j)):0\le j\le8\},\\
 \partial A_{120}&=\{((i,-1),(i,0)),((i,8),(i,9)):0\le i\le8\}.
\end{align*}
Each collection contains eighteen links. Their labels retain channel
provenance; each link carries a number operator with spectrum
$\{0,1,2\}$. Let $N_{90},N_{120}$ be the corresponding sums.

The angular areal scale retains the counterangle already constructed:
\[
 \theta_{\rm clk}=\frac{\pi C^*_{\deg}}{1080},\qquad
 a_0=\frac{1+2\cos(\pi/18)}3\,\theta_{\rm clk}.
\]
The prefactor is the normalized real trace of
$\operatorname{diag}(1,e^{i\pi/18},e^{-i\pi/18})$.
Thus $a_0$ retains its meaning as a dimensionless area normalization.
The area readout is taken on the positive branch
$\theta_{\rm clk}>0$ and $\gamma>0$.
With the Barbero--Immirzi functional $\gamma$ and the length already
determined,
\[
 \mathcal A=\gamma a_0L_\alpha^2(N_{90}+N_{120}).
\]
The thirty-six local operators act on distinct tensor factors.
The spectrum of $N=N_{90}+N_{120}$ is therefore
$\{0,\ldots,72\}$, and the multiplicity of $n$ is the coefficient
of $z^n$ in $(1+z+z^2)^{36}$. The proof consists in multiplying
the generating functions of the factors. The new length transports
the areal unit and preserves this combinatorics.

\begin{proposition}[Joint readout of area and sectoral provenance]
The pair
$N=N_{90}+N_{120}$ and $\mathcal D_{\partial}=90N_{90}+120N_{120}$
recovers the two sectors:
\[
 N_{90}=4N-\frac{\mathcal D_{\partial}}{30},\qquad
 N_{120}=\frac{\mathcal D_{\partial}}{30}-3N.
\]
\end{proposition}
\begin{proof}
The transformation matrix is
$\left(\begin{smallmatrix}1&1\\90&120\end{smallmatrix}\right)$,
with determinant $30$. Its inverse gives the formulas. Commutativity
of the two number operators implies that of the two readout operators.
\end{proof}
In the modular realization with
$\rho_A=Z_A^{-1}\exp(-\Delta_{\rm mod}\mathcal D_{\partial})$,
$\Delta_{\rm mod}>0$, functional calculus gives
$-\log\rho_A=(\log Z_A)I+\Delta_{\rm mod}\mathcal D_{\partial}$.
With the normalizations retained, the pair consisting of area and
the modular Hamiltonian recovers the two occupations. The reconstruction
proof holds for each constructed value of $\Delta_{\rm mod}$; the new
longitudinal readout operator does not redefine that coefficient or
the modular state. The sum of occupations and sectoral provenance
express distinct observables. Likewise, the thirty-six response links
and the eighty-one unit capacities of the cubic cut retain their roles.

\subsection{Joint variation and compatible memory reduction}
\label{md:rigidez:variaciones}

For norm-differentiable collections that are uniformly positive in a
neighbourhood, let $\kappa=L_\alpha^2/(\hbar c)$. Then
\[
 \delta R=\delta\kappa\,H+\kappa\,\delta H,\qquad
 \frac{\delta G}{G}=3\frac{\delta c}{c}
 +2\frac{\delta L_\alpha}{L_\alpha}-\frac{\delta\hbar}{\hbar}.
\]
The first identity follows by differentiating
$R=L_\alpha Y$ and $H=(\hbar c/L_\alpha)Y$; it does not require
$[Y,\delta Y]=0$. An admissible connection variation acting on $Y$
and preserving the three sections maintains $\delta G=0$.
The statement covers variations of this radial realization;
equivalence between complete field functionals has its own domains
and boundary conditions.

Now let
$Y=\left(\begin{smallmatrix}A&C\\C^*&D_i\end{smallmatrix}\right)>0$
under the separation of readout and memory, with
$S=A-CD_i^{-1}C^*>0$. The minimizing extension of a vector $b$ is
$(b,-D_i^{-1}C^*b)$, as completing the square shows:
\[
 \left\langle\binom b m,Y\binom b m\right\rangle
 =\langle b,Sb\rangle+
 \langle m+D_i^{-1}C^*b,D_i(m+D_i^{-1}C^*b)\rangle.
\]
The coherent reduction therefore gives
\[
 H_{\rm eff}=\frac{\hbar c}{L_\alpha}S,\qquad
 R_{\rm eff}=L_\alpha S,\qquad\Lambda_b=L_\alpha S^{-1}.
\]
The compression of $Y^{-1}$ onto the readout block is $S^{-1}$;
the direct compression of $Y$ is $A$. Memory distinguishes these operations.

If dynamical elimination produces a bounded temporal metric
$Z\ge z_0I$, $z_0>0$, take $T=Z^{1/2}$, bounded and invertible.
Canonical normalization requires the dual transports
\[
 H_c=T^{-*}H_{\rm eff}T^{-1},\qquad
 R_c=T^{-*}R_{\rm eff}T^{-1},\qquad
 \Lambda_c=T\Lambda_bT^*.
\]
Multiplying them yields $H_c\Lambda_c=\hbar cI$ and
$R_c\Lambda_c=L_\alpha^2I$, without a commutativity hypothesis.
In the resolvent expansion of an energy with blocks $H_{bb},V,H_{ii}$,
with $H_{ii}$ positive and invertible and
$|\hbar\omega|<\|H_{ii}^{-1}\|^{-1}$,
\[
 K_b(\omega)=H_{bb}-\hbar\omega I
 -V(H_{ii}-\hbar\omega I)^{-1}V^*
 =H_{\rm stat}-\hbar\omega Z+O(\omega^2),
 \qquad Z=I+VH_{ii}^{-2}V^*.
\]
The last expression is local around zero and retains the order of
approximation; the preceding dual products are exact identities
of the normalized readout operators.

Proportionality is also preserved under full dynamical elimination.
For $z$ such that $H_{ii}-zI$ is invertible, define
\[
 \mathcal K_H(z)=H_{bb}-zI-V(H_{ii}-zI)^{-1}V^*,\qquad
 \kappa=\frac{L_\alpha^2}{\hbar c}.
\]
The same readout--memory decomposition applied to $R=\kappa H$
produces
\begin{equation}
 \mathcal K_R(\kappa z)=\kappa\mathcal K_H(z).
 \label{md:rigidez:resolvente-exacto}
\end{equation}
Indeed, its blocks are $R_{bb}=\kappa H_{bb}$,
$R_{bi}=\kappa V$ and $R_{ii}=\kappa H_{ii}$, and
\[
 (R_{ii}-\kappa zI)^{-1}
 =\kappa^{-1}(H_{ii}-zI)^{-1}.
\]
Substitution into the Schur complement proves
\eqref{md:rigidez:resolvente-exacto}, preserving the order of the factors.
When both complements are invertible, their inverses satisfy
$\mathcal K_R(\kappa z)^{-1}=\kappa^{-1}\mathcal K_H(z)^{-1}$.
The energy parameter is $z=\hbar\omega$, and its image in the radial
readout operator is $\kappa z$: both retain their respective dimensions.
The identity holds throughout this resolvent domain and therefore
jointly transports all orders of dynamical memory. The preceding local
expansion describes its first coefficients; this equality explains
the exact preservation of radial proportionality under reduction.

\subsection{Temporal rigidity and boundary memory}
\label{md:rigidez:memoria}

Preserving only area allows $R\mapsto\zeta R$,
$\Lambda\mapsto\Lambda/\zeta$. Preserving $H,\hbar,c$ and
$H\Lambda=\hbar cI$ as well forces $\zeta=1$. If
$H\mapsto\zeta H$ is also applied, the quotient $R/H$ and the
coefficient $G$ remain unchanged. Comparisons must preserve the same
joint observations.

The lifted phase supplies another effective constraint. In a bounded
realization $H=H_{\rm clk}+D_0^*WD_0$ with $W>0$, the evolution
$P(t)=\exp(-itH/\hbar)$ determines $H=i\hbar\dot P(0)$.
With $H_{\rm clk},D_0,\hbar$ and the clock fixed, equality of the
evolutions implies $D_0^*(W'-W)D_0=0$. For a route with fixed
initial endpoint, $D_0$ is triangular and invertible, so $W'=W$ follows.
For a general graph, the form is determined on $\operatorname{Ran}D_0$;
subsequently extending the residue domain would change the elimination
problem. A single return $P(\tau)$ admits spectral aliases, whereas
the lifted phase or compatible times tending to zero distinguish them.

In the nonadic capacity chart, let
$\rho=\log_{10}9$, $\delta=1-\rho$,
$p_j=\lfloor j/\delta\rfloor$ and $v_j=p_j-j$. The capacity of a window is
\[
 C(j,n)=p_{j+n}-p_j-n.
\]
For $n=35$ its values are $729$ or $730$, depending on the origin;
capacity return modulo nine distinguishes the former within this class.
The window $j=20\to55$ satisfies $p_j:437\to1201$ and
$v_j:417\to1146$. Its capacity phase returns to $3$, and the quotient
increases from $46$ to $127$. The six capacities of $20$ and the
twenty-nine of $21$ sum to $729$; accumulated memory remains part
of the state.

For the first arrival $T(k)=\lceil k/\rho\rceil-1$, the count is
$N(k,C)=T(k+C)-T(k)-C$. Writing
$\epsilon(k)=\lceil k/\rho\rceil-k/\rho$ gives exactly
\[
 N(k,C)=(\rho^{-1}-1)C+\epsilon(k+C)-\epsilon(k).
\]
Capacity-$729$ windows with counts $34$ and $35$ retain different
boundary terms. Even $k=21$ and $k=2289$ share
$(k\bmod108,T(k)\bmod108)=(21,22)$ and accumulate $1$ and
$109$ vacancies, respectively. Coincidence of finite calendars retains
less information than the complete state. These calculations interpret
the return; the longitudinal readout operator in
\eqref{md:rigidez:longitud-def} has its own construction map and does
not select its scale through these windows.

\subsection{Geometric normalization and evaluation}
\label{md:rigidez:evaluacion}

The central APP block
$B_c=\left(\begin{smallmatrix}7&2\\2&7\end{smallmatrix}\right)$
satisfies $B_c^{\mathsf T}g_cB_c=45g_c$,
$g_c=\operatorname{diag}(1,-1)$. Its positive unimodular normalization is
$B_c/\sqrt{45}$: an additional positive factor $a$ would give determinant
$a^2$. The areal--volumetric incidence
$T_{\rm av}=\left(\begin{smallmatrix}1&1\\1&2\end{smallmatrix}\right)$
preserves
$G_{\rm av}=\left(\begin{smallmatrix}2&1\\1&-2\end{smallmatrix}\right)$
by direct multiplication.
In a Minkowski realization, the causal form and dimensional soldering
have distinct roles. The cochain
$\beta_m=\ell_m\sigma_mW_\square e_\lambda$,
$\ell_m=\ell_0/9^m$, and the relation $\ell_0=ct_0$ retain the scale
together with the cone. Altering the soldering while preserving only
the signature changes that realization.

The expanding Pell unit with norm one and trace four satisfies
$x+x^{-1}=4$, $x>1$, and hence $x=2+\sqrt3$.
The orbit $r_n=r_0x^n$ preserves $r_nr_{-n}=r_0^2$.
When $r_0=L_\alpha$, the product preserves the same areal section.
The norm, marked form, return, and longitudinal section constrain
different objects; their roles are retained when they are composed.

\Needspace{10\baselineskip}
In the declared SI section, $L_*=1\,\mathrm m$,
$S_*=10^{-34}\,\mathrm{J\,s}$ and $c=299792458\,\mathrm{m\,s^{-1}}$,
structural substitution produces
\[
 L_\alpha=1.616254982625126330\ldots\,10^{-35}\,\mathrm m,\qquad
 \hbar_{\rm ret}=1.054571816915039061\ldots\,10^{-34}\,\mathrm{J\,s},
\]
\begin{equation}
 G_{\rm ret}=6.674299659414022706\ldots\,10^{-11}
 \,\mathrm{m^3\,kg^{-1}\,s^{-2}}.
 \label{md:rigidez:valor}
\end{equation}
The evaluator substitutes the regional coordinates and register $K$,
reconstructs $\alpha$, the return of $H_5$, and $L_\alpha$;
metrological comparison follows. The CODATA 2022 reference is
$6.67430(15)\,10^{-11}$ in those units \cite{codata2022}.
The difference is $-3.40586\,10^{-18}$, equivalent to
$-0.00227$ standard uncertainties. The calculated digits describe the
mathematical evaluation in this section, not an experimental uncertainty.
The pre-return and post-return action sections retain their labels and
$\hbar_aG_a=c^3L_\alpha^2$.

The structural meaning of $G$ is thereby determined: it is the
coefficient that converts the energy readout into conjugate length in
the common radial realization while simultaneously preserving action
and area. Memory, sheet markers, and bases belong to the construction
that fixes this coefficient. Their joint preservation explains the
proved uniqueness.
